% Arquitectura del clasificador de tipo de golpe (sección «Tipo de golpe»), al
% estilo del recuadro BST-0 de la figura 2 de BST, con los colores de la memoria
% y los iconos del jugador y la pelota en lugar de cajas para las entradas (con
% sus dimensiones, sin nombre: el icono ya dice qué es). Las dimensiones van con
% la notación del texto y de las ecuaciones: T pasos, J puntos clave, D dimensión
% del modelo y C clases.
% Sigue a ShotTypeBST (padel_ml.shot_type_model): un TCN por secuencia, un
% codificador temporal con los mismos pesos para las dos, la atención cruzada
% (Q de la pose; K y V de la pelota), el codificador de interacción y la cabeza
% sobre los tres resúmenes. Colores, estilos e iconos en config/config.tex.
\begin{adjustbox}{max width=\linewidth}
\begin{tikzpicture}[
  op/.style={circle, draw=tinta2, fill=white, line width=0.5pt, minimum size=3.6mm,
    inner sep=0pt, font=\scriptsize, text=tinta},
  token/.style={draw=tinta2, fill=white, line width=0.5pt, minimum size=2.4mm, inner sep=0pt},
  rama/.style={circle, fill=tinta2, minimum size=1.3mm, inner sep=0pt},
  alta/.style={caja, minimum height=31mm},
  modulo/.style={caja, text width=10mm},
  nombre/.style={font=\footnotesize, text=tinta, inner sep=1pt, align=center},
  pics/posicion/.style={code={
    \node[op] (-c) {};
    \draw[tinta2, line width=0.45pt]
      plot[domain=-0.12:0.12, samples=25, smooth] (\x, {0.065*sin(\x/0.12*180)});
  }},
]
  \def\yP{1.15}
  \def\yB{-1.15}
  \def\yR{2.15}

  % Entradas: el jugador que golpea y la pelota
  \pic[scale=1.25] at (0,\yP-0.12) {jugador=tinta};
  \node[nombre, below, font=\scriptsize, text=tenue] at (0,\yP-0.16) {$T\times J\times3$};
  \pic[scale=1.7] at (0,\yB+0.2) {pelota=pelota};
  \node[nombre, below, font=\scriptsize, text=tenue] at (0,\yB-0.04) {$T\times3$};
  \node[modulo] (tcnP) at (1.75,\yP) {TCN};
  \node[modulo] (tcnB) at (1.75,\yB) {TCN};
  \draw[flecha] (0.55,\yP) -- (tcnP);
  \draw[flecha] (0.55,\yB) -- (tcnB);

  % Token de clase y codificación posicional antes del codificador temporal
  \foreach \r/\y in {P/\yP, B/\yB} {
    \node[op] (cat\r) at (2.85,\y) {$\|$};
    \node[token] (tok\r) at (2.85,\y+0.5) {};
    \node[op] (sum\r) at (3.4,\y) {$+$};
    \pic (pe\r) at (3.4,\y-0.5) {posicion};
    \draw[flecha] (tcn\r) -- (cat\r);
    \draw[flecha] (tok\r) -- (cat\r);
    \draw[flecha] (cat\r) -- (sum\r);
    \draw[flecha] (pe\r-c) -- (sum\r);
  }

  \node[alta, text width=17mm] (enc1) at (4.75,0)
    {Codificador\\temporal\subcaja{2 capas}\subcaja{mismos pesos}};
  \draw[flecha] (sumP) -- (sumP -| enc1.west);
  \draw[flecha] (sumB) -- (sumB -| enc1.west);

  % Se separa el resumen (token) de la secuencia, que sigue
  \node[rama] (sP) at (6.0,\yP) {};
  \node[rama] (sB) at (6.0,\yB) {};
  \draw[linea] (enc1.east |- sP) -- (sP);
  \draw[linea] (enc1.east |- sB) -- (sB);

  \node[alta, text width=15mm] (cross) at (8.6,0) {Atención\\cruzada\subcaja{1 capa}};
  \node[op] (sum2P) at (7.15,\yP) {$+$};
  \pic (pe2P) at (7.15,\yP-0.5) {posicion};
  \node[op] (sum2B) at (7.15,\yB) {$+$};
  \pic (pe2B) at (7.15,\yB-0.5) {posicion};
  \draw[flecha] (sP) -- (sum2P) node[pos=0.5, above, font=\scriptsize, text=tenue] {$T\times D$};
  \draw[flecha] (sB) -- (sum2B) node[pos=0.5, above, font=\scriptsize, text=tenue] {$T\times D$};
  \draw[flecha] (pe2P-c) -- (sum2P);
  \draw[flecha] (pe2B-c) -- (sum2B);
  \draw[flecha] (sum2P) -- (sum2P -| cross.west) node[pos=0.5, above, font=\scriptsize] {$Q$};
  \draw[flecha] (sum2B) -- (sum2B -| cross.west) node[pos=0.5, above, font=\scriptsize] {$K, V$};

  % Codificador de interacción sobre la secuencia fusionada
  \node[op] (catF) at (10.0,0) {$\|$};
  \node[token] (tokF) at (10.0,0.5) {};
  \node[op] (sumF) at (10.55,0) {$+$};
  \pic (peF) at (10.55,-0.5) {posicion};
  \node[caja, text width=21mm] (enc2) at (12.15,0) {Codificador\\de interacción\subcaja{1 capa}};
  \draw[flecha] (cross) -- (catF);
  \draw[flecha] (tokF) -- (catF);
  \draw[flecha] (catF) -- (sumF);
  \draw[flecha] (peF-c) -- (sumF);
  \draw[flecha] (sumF) -- (enc2);

  % Cabeza: los tres resúmenes, concatenados
  \node[op] (head) at (13.75,0) {$\|$};
  \node[modulo] (mlp) at (14.75,0) {MLP};
  \node[entrada, text width=12mm] (out) at (16.2,0) {$C$ valores\subcaja{uno por clase}};
  \draw[flecha] (sP) -- (6.0,\yR) -| (head);
  \draw[flecha] (sB) -- (6.0,-\yR) -| (head);
  \pic[scale=0.8] at (9.6,\yR+0.08) {jugador=tinta};
  \node[dato, anchor=west] (rP) at (9.8,\yR+0.18) {resumen de la pose};
  \pic[scale=1.1] at (9.6,-\yR-0.2) {pelota=pelota};
  \node[dato, anchor=west] (rB) at (9.8,-\yR-0.2) {resumen de la pelota};
  \draw[flecha] (enc2) -- (head);
  \draw[flecha] (head) -- (mlp);
  \draw[flecha] (mlp) -- (out);

  \begin{scope}[on background layer]
    \node[panel, fill=panelgolpe, fit={(-0.75,0) (out) (rP) (rB) (6.0,2.75) (6.0,-2.6)}] (g) {};
  \end{scope}

  % Leyenda, centrada bajo el panel
  \begin{scope}[shift={($(g.south)+(-7.75,-0.45)$)},
      leyenda/.style={font=\scriptsize, text=tinta2, anchor=west, inner sep=1pt}]
    \node[token] (l1) at (0,0) {};
    \node[leyenda, right=1mm of l1] (t1) {token de clase};
    \node[op, right=5mm of t1] (l2) {$\|$};
    \node[leyenda, right=1mm of l2] (t2) {concatenación};
    \node[op, right=5mm of t2] (l3) {$+$};
    \node[leyenda, right=1mm of l3] (t3) {suma};
    \path (t3.east) ++(6.8mm,0) pic (l4) {posicion};
    \node[leyenda, right=1mm of l4-c] (t4) {codificación posicional};
    \node[rama, right=5mm of t4] (l5) {};
    \node[leyenda, right=1mm of l5] {se separa el resumen de la secuencia};
    \node[leyenda] at (-0.12,-0.45)
      {$T$: pasos \quad $J$: puntos clave \quad $D$: dimensión del modelo \quad $C$: clases};
  \end{scope}
\end{tikzpicture}
\end{adjustbox}
